% APÉNDICE I

\VIUSectionStar{Apéndice I}

Este apéndice recoge información complementaria de autoría propia relacionada con la operación del prototipo implementado y con la organización técnica del artefacto. Su finalidad es ampliar, sin interrumpir la lectura del cuerpo principal, aspectos que ayudan a interpretar el estado actual del módulo y sus condiciones de ejecución.

En particular, el prototipo \textit{AQRisk} dispone de una arquitectura de ejecución compuesta por un núcleo Python, una API HTTP, una interfaz web construida para control y explicabilidad, y un entorno de despliegue mediante Docker. Esta estructura permitió verificar el flujo completo del sistema en condiciones reproducibles y preparar la captura de evidencia visual para la memoria.

Desde el punto de vista documental, el apéndice sirve para ubicar elementos de apoyo que complementan el análisis principal, mientras que los anexos reúnen tablas extensas y material técnico de consulta. En esta versión del informe, los anexos incorporan el checklist de calidad, el diccionario operativo de extracción, la evidencia cuantitativa complementaria y la base completa de reglas del motor difuso.
